\answer{ex:09:1}{(a)~$P_\infty=0.2$, memoria de $20\,\text{s}$. (b)~$P_\infty=1$, memoria de $1\,\text{s}$. (c)~La memoria es proporcional a la amplitud del ruido: seis veces.}
\answer{ex:09:2}{$P(t)=P_\infty\coth(t\sqrt{q/R})$. (a)~$\frac12\sqrt{R/q}=5\,\text{s}$, la mitad que la memoria. (b)~$t=\frac12\ln21\,\sqrt{R/q}\approx15\,\text{s}$: lo que debe mantenerse elevado \nt{ACET}.}
\answer{ex:09:3}{Varianza $qT/2+R/(2T)$; $T^*=\sqrt{R/q}$, varianza $\sqrt{qR}=P_\infty$, como en el filtro~\eqref{eq:kalman}; crece $(\kappa+1/\kappa)/2$ veces: $1.25$ y $5.05$.}
\answer{ex:09:4}{$a>a_w=1-\theta/(mw)$: $0.15$ con $w=0.041$, $0.22$ con $w=0.045$.}
\answer{ex:09:5}{Los errores aparecen con $M\approx40$--$60$; con $M=80$ sobran $1.9$ neuronas; con $M=100$, fracción del patrón $0.86$; la interferencia de los demás patrones tiene varianza $\approx(M-1)p(1-p)/N$, así que $M_{\max}\approx80$, $M_{\max}\propto N/p$.}
\answer{ex:09:6}{Por ejemplo, $\eta(a)=\eta\min\bigl(1,\max(0,(a-0.3)/0.6)\bigr)$. Con $\eta a$ las tres neuronas ajenas entraron en el patrón; con $\eta(a)$, ninguna, y la escritura con $a\ge0.9$ es la misma (fracción $1.00$).}
\answer{ex:09:7}{(a)~$40$ repeticiones; con la $\eta(a)$ del ejercicio~\ref{ex:09:6}, nunca. (b)~$200$ repeticiones, $10$ noches.}
